\begin{abstract}
How much of Colombia's measured output is produced by informal work? The answer cannot be read from the labor-informality rate: Colombia combines an informality rate above 50\% with a much smaller informal share of output, because informal work is less productive. We estimate a Bayesian two-sector business-cycle model on quarterly national accounts and on reconstructed employment and informality series for 1993--2026. Identification comes from a new quarterly series of informal-to-formal labor earnings built from household-survey microdata for 2007--2026. Formal production uses capital and labor, informal production is labor based, and an accounting relationship maps relative earnings, the informality rate and the formal capital share into the informal output share. We show that aggregate data alone cannot pin down the size of the informal sector, and that relative earnings identify its quarterly path: from 2007 the path is invariant, within 0.3 percentage points, to the model's steady state. The informal share declines from about 30\% of GDP in 1993 to 20\% in 2026Q2, while labor informality falls from about 69\% to 54\%. The estimates lie within the range of published shadow-economy series for Colombia, close to the enhanced MIMIC estimates of Asllani, Dell'Anno and Schneider for 2009--2020. Sensitivity analysis quantifies the role of the capital share and of non-wage labor costs.
\end{abstract}

\noindent\textbf{Keywords:} informal output; labor informality; Bayesian DSGE; identification; relative earnings; Colombia

\noindent\textbf{JEL:} E26, E32, C11, O17
